\documentclass[11pt]{article}
\usepackage[T1]{fontenc}
\usepackage{textcomp}
\usepackage[margin=0.75in]{geometry}
\usepackage{amsmath,amssymb,amsthm}
\usepackage{array,booktabs}
\usepackage{hyperref}
\setlength{\parindent}{0pt}
\newtheorem{theorem}{Theorem}
\theoremstyle{definition}
\newtheorem{definition}{Definition}
\title{Flow Proof Artifact: Ring-one-right-identity}
\author{Generated by \texttt{flow doc proof}}
\date{Flow Proof Book}
\begin{document}
\maketitle
One is the right multiplicative identity in a ring.\\[0.5em]
\textbf{Source.} Ring one right identity\\[0.5em]
\bigskip
\noindent{\large \textbf{Derived fact 1.} \textit{One is the right multiplicative identity in a ring} \hfill Ring \cdot one \cdot right identity}\par
\smallskip\noindent\textit{Coordinate: Ring \cdot one \cdot right identity}\par
\smallskip\noindent\textit{Source: Ring one right identity}\par
\smallskip\noindent\textit{Built on: \hyperref[thm:Ring-one-right_identity]{Derived fact 1}}\par
\medskip
\noindent\textbf{Goal.} One is the right multiplicative identity in a ring\par
\noindent$\displaystyle \forall x \in \guillemotleft{}Ring\guillemotright{} \guillemotleft{}Nat\guillemotright{}\quad «Ring» «mul»(x, 1) = x$\par
\medskip
\noindent
\begin{tabular}{@{} >{\raggedright\arraybackslash}p{0.06\textwidth} >{\raggedright\arraybackslash}p{0.47\textwidth} >{\raggedleft\arraybackslash}p{0.41\textwidth} @{}}
\textbf{\#} & \textbf{Proof} & \textbf{Mathematics} \\
\midrule
\textbf{1.} & We prove that right identity for one on Ring. & \\[0.45em]
\textbf{2.} & We invoke the derived fact: \guillemotleft{}Ring\guillemotright{} \guillemotleft{}mul\guillemotright{} (instantiated for x, 1). & \\[0.45em]
\textbf{3.} & From step 2, this implies \guillemotleft{}Ring\guillemotright{} \guillemotleft{}mul\guillemotright{}(x, 1) equals x. Hence proven. & $\displaystyle «Ring» «mul»(x, 1) = x$ \\[0.45em]
\bottomrule
\end{tabular}
\label{thm:Ring-one-right_identity}
\par\medskip\hrule
\end{document}
